% !TeX root = ../../main.tex
\documentclass[../../main.tex]{subfiles}
\begin{document}

\paragraph{Literature on gas jet configuration}

\todo{$\succ$ Add literature to water atomizer project}

Motaman et al. \cite{Motaman2013NumericalExperimental} investigate numerically and experimentally how the gas jet mismatch angle and stand-off distance affect back-stream flow in close-coupled gas atomization, showing that smaller mismatch angles and lower gas pressures promote back-stream, while larger angles ($\geq 5^\circ$) or higher pressures suppress or eliminate it. This paper reminds the pros and cons of the free-fall and close-coupled gas atomization, where the gas-to-melt interaction is easier to control in free-fall gas atomization, while the particle size distribution and the process efficiency is greater in close-coupled gas atomization. Indeed, either the melt stream and gas jet are easier to control in close-coupled gas atomization, and less energy dissipation occurs leading to reduced gas consumption. The experiment have been done with single stage discrete jet nozzles with convergent-divergent dies. Using a regular dies (cylindrical discrete jets as an example) leads to uncontrolled expansion of the jet creating some shocks and reducing the velocity. In contrast, a convergent-divergent dies can be designed so as to have the environment pressure at the nozzle outlet. In the presented configuration, some backflow is experienced, resulting to metal freeze at the gaz nozzle tips. 

\paragraph{Set-up description}

\subparagraph{Hardware and nozzle}

\begin{itemize}
    \item Bosch aluminum profiles frame and fittings
    \item FDM printed gas spreader for stage 1 and 2
    \item Festo connectors and tubings
    \item FDM printed customized nozzles
\end{itemize}

\subparagraph{Sensor and electronics}

\begin{itemize}
    \item Pressure transducer -1 to 10 bar
    \item National Instrument input module
\end{itemize}

\paragraph{Experimental results}

\begin{itemize}
    \item Start simple with single stage
    \text{P\_!\_alpha\_0\_....}
\end{itemize}

\subparagraph{Single-stage nozzle}

\subparagraph{Dual-stage nozzle}

\paragraph{Numerical simulation}

\subparagraph{Literature on gas atomization flow dynamics}

\subparagraph{Model description and geometrical assumptions}

\subparagraph{Mesh study and computation method}

\subparagraph{Boundary conditions}

\subparagraph{Computational results}









\paragraph{Pre-experiment ideas}

\begin{itemize}
    \item Vertical water drop on horizontal air jet, then try different angle of air jet until vertical.\\ When vertical (water and air axis are parallel then), we should see the suction effect: the drop wants to go in the jet.
    \item From that idea, what if multiple jets are directed vertically next to the drop axis (distance?) just to delaminate the water stream?\\
\end{itemize}

\paragraph{Nozzle design ideas}

\vspace{0mm}
\begin{minipage}{\linewidth}
    FIRST STAGE
        \begin{itemize}
        \item[$\Rightarrow$] focus metal stream
        \item[$\Rightarrow$] protect from backflow
        \item[$\Rightarrow$] create a vacuum effect and first slight breakup
        \end{itemize}
            \begin{itemize}
            \item minimum apex angle
            \item annular (and swirling, to decrease even more apex angle)
            \item low pressure
            \item starting before the melt nozzle end, guided but the melt nozzle. (In the future, the melt nozzle and the gaz nozzle should be integrated in one nozzle, in printed ceramic)
            \end{itemize}
\end{minipage}
\vspace{5mm}


\vspace{0mm}
\begin{minipage}{\linewidth}
SECOND STAGE
        \begin{itemize}
        \item[$\Rightarrow$] suck the 1st stage flow and the liquid stream
        \item[$\Rightarrow$] secondary breakup
        \end{itemize}
            \begin{itemize}
            \item very small apex angle
            \item Laval geometry
            \item annular and swirling
            \item high pressure, just above the closed-wake pressure
            \end{itemize}
\end{minipage}
\vspace{0mm}


\paragraph{Printing tests for nozzles}

\begin{itemize}
    \item stage1, only 3 slots with pressure transducer slot, side wall 1.5 mm: test sealing with one PU inlet and two clogged ones, outlet clogged.
    \item stage1, only 3 slots with pressure transducer slot, side wall 1.5 mm: test sealing with one PU inlet and two clogged ones, outlet open.
    \item Stage1, full manifold: test sealing with one PU inlet and all the outlet closed.
\end{itemize}

Results: as expected, small wall thickness brings to more porosity. Choice to go with full of plastic, however, under high pressure, still porous. In the real life, the pressure 
inside the manifold will be lower since the circuit will be open.

\paragraph{Parameters and configuations}

\todo{$\succ$ Make a scheme with all the parameters}

\vspace{0mm}
\begin{table}[H]
    \centering
    \begin{tabular}{l c c}
        \toprule
        \textbf{} & \textbf{} & \textbf{}  \\
        \midrule
        Jet type & Discrete & Annular slit   \\
        Jet geometry & Regular / Laval & Regular / Laval   \\
        Apex angle & From ?? to ?? & From ?? to ??   \\
        Swirl angle & From ?? to ?? & From ?? to ?? \\
        D1 &From ?? to ?? & From ?? to ?? \\
        D2 & From ?? to ?? & From ?? to ??  \\
        Pressure & From ?? to ?? & From ?? to ??  \\
        \bottomrule
    \end{tabular}
    \caption{Parameters of the water atomizer nozzles}
    \label{}
\end{table}
\FloatBarrier
\vspace{0mm}

report, surf integral, area weighted av, fiels var\\
keep the inlet pressure constant.

\paragraph{Experimental Set-up}

\subparagraph{Instrumentation}

\begin{itemize}
    \item Pressure transducer: Endress+Hauser Cerabar PMC11, 4 to 20V, -1 to 4 bar
    \item Current Input Module: NI 9203 with its corresponding chassis
\end{itemize}

\paragraph{Water inlet}

Flow rate and fowl speed estimation for the water nozzle:\\

\vspace{0mm}
\begin{equation}
\begin{aligned}
Q = C \cdot A \cdot \nu
\hspace{10mm}
with
\hspace{10mm}
\nu = \sqrt{2 \cdot g \cdot h}
\label{Flow_rate}
\end{aligned}
\end{equation}
\vspace{0mm}

\vspace{0mm}
\begin{minipage}{\linewidth}
Where:
\begin{itemize}
   \item $Q$ [m³/s]: flow rate
   \item $C$ [-]: discharge coefficient
   \item $A$ [m²]: area
   \item $\nu$ [m/s]: velocity
   \item $g$ [m/s²]: gravity acceleration
\end{itemize}
\end{minipage}
\vspace{0mm}

By taking:
\begin{itemize}
    \item C = 0.6,
    \item water nozzle diameter = 4~mm,
    \item h = 2/3 * 25 cm = 0.017 m (equivalent value for a 2L Coca-Cola bottle),
    \item g = 9.81 m/s²
\end{itemize}

$\mu$ $\approx$ 1.81 m/s
Q $\approx$ 0.0137 L/s $\Rightarrow$ $\approx$ 0.83 L/min $\Rightarrow$ 2~min~25~s

\todo{$\succ$ just air in water nozzle?}
\todo{$\succ$ air in just 2 jets of discrete jets?}
\todo{$\succ$ make a list of the different configurations to test}
\todo{$\succ$ Calculation model parametric Fluent}
... with vertical and horizontal probes lines\\


\end{document}
